\providecommand{\pdfxopt}{a-1b,mathxmp}

%\documentclass[10pt]{article}
%\usepackage{beamerarticle}
\documentclass[ignorenonframetext,aspectratio=169]{beamer}

\input{/home/asreeve/current/tex_style/preamble_102}

\setbeamertemplate{enumerate item}[square]
\setbeamertemplate{enumerate subitem}[circle]
\setbeamertemplate{enumerate subsubitem}[square]
\setbeamercolor*{enumerate item}{fg=red,bg=blue}
\setbeamercolor*{enumerate subitem}{fg=blue,bg=red}

\begin{document}
\begin{frame}{A $CO_2$ 'box model'}
  \begin{minipage}{.58\textwidth}
    The diagram below illustrates an idealized model of $CO_2$ in the atmosphere. $CO_2$ fluxes into and out of the atmosphere are lumped into two processes: emissions to the atmosphere and residence time of $CO_2$ (how quickly $CO_2$ is lost).    
    
    \includegraphics[width=\textwidth]{../../figs/CO2_Model.pdf}
  \end{minipage} \quad
  \begin{minipage}{.38\textwidth}
    Four parameters are used to simulate $CO_2$ concentrations ($CO_2$):
    \begin{relsize}{-1}
    \begin{itemize}
    \item The starting conc. of $CO_2$ (C0)
    \item Initial emission rate (Em)
    \item Growth (change) in emission rate over time (GR)
    \item $CO_2$ Residence time (RT) 
    \end{itemize}
  \end{relsize}
  Mathematically, changes in $CO_2$ Conc. with time (t) are:
    \begin{equation*}
      \Delta CO_2 = \Delta t \left(Em - \frac{CO_2}{RT} \right)
    \end{equation*}
  \end{minipage}
\end{frame}

\begin{frame}{Spreadsheet Structure}
\end{frame}

\begin{frame}{Exploring Model Behavior}
\end{frame}

\begin{frame}{Model Calibration}
\end{frame}

\begin{frame}{Making Predicitions}
\end{frame}

\begin{frame}{Relating $CO_2$ Conc. to Temperature}
\end{frame}

\end{document}
